\chapter{线性微分方程}

在这一章中，我们将研究在复平面中的二阶微分方程

\begin{equation*}
\frac {d ^ {2} f}{d z ^ {2}} + p (z) \frac {d f}{d z} + q (z) f = 0.\tag{7.1}
\end{equation*}

显然贝塞耳方程和埃尔米特方程为其特殊情形。方程(7.1)的性质，由函数 $p(z)$ 和 $q(z)$ 所确定。

如果在点 $z = z_0$ 处函数 $q(z)$ 和 $p(z)$ 解析，则 $z = z_0$ 称为方程的正常点。在其余情形， $z = z_0$ 称为奇点。

\section{在正常点附近方程的解}

假定 $z_0$ 是正常点，并设 $z = z_0 + \xi$ （这里 $|\xi| \leqslant a$ ）。还假定， $p(z)$ 和 $q(z)$ 在以 $z_0$ 为中心， $a$ 为半径的圆内（也在圆上）解析。为以后方便计，作未知函数 $f(z)$ 的变换，设

\begin{equation*}
f (z) = e ^ {- \frac {1}{2} \int_ {z _ {0}} ^ {z} p (z) d z} g (z),
\end{equation*}

则
\begin{equation*}
\frac {d f}{d z} = \left[ - \frac {1}{2} p (z) g (z) + g ^ {\prime} (z) \right] e ^ {- \frac {1}{2} \int_ {z _ {0}} ^ {z} p (z) d z},
\end{equation*}

\begin{equation*}
\frac {d ^ {2} f}{d z ^ {2}} = \left[ - \frac {1}{2} p ^ {\prime} (z) g (z) - p (z) g ^ {\prime} (z) + g ^ {\prime \prime} (z) + \right.
\end{equation*}

\begin{equation*}
\left. + \frac {1}{4} p ^ {2} (z) g (z) \right] e ^ {- \frac {1}{2} \int_ {z _ {0}} ^ {z} p (z) d z}.\tag{7.2}
\end{equation*}

将(7.2)代入(7.1)就得关系函数 $g(z)$ 的方程：

\begin{equation*}
\frac {d ^ {2} g (z)}{d z ^ {2}} + \left[ q (z) - \frac {1}{4} p ^ {2} (z) - \frac {1}{2} \frac {d p (z)}{d z} \right] g (z) = 0.\tag{7.3}
\end{equation*}

在函数 $q(z)$ 和 $p(z)$ 解析的区域内，变换(7.2)将 $f(z)$ 解析的区域变为 $g(z)$ 解析的区域，因此新方程的正常点与原方程的正常点一致。因 $z_0$ 固定，故可选取 $\xi$ 作为新的自变量：

\begin{equation*}
\frac {d ^ {2} g (\xi)}{d \xi^ {2}} + v (\xi) g (\xi) = 0,\tag{7.4}
\end{equation*}

这里
\begin{equation*}
v (\xi) = q (\xi) - \frac {1}{4} p ^ {2} (\xi) - \frac {1}{2} p ^ {\prime} (\xi).
\end{equation*}

为了求解方程(7.4)，可考虑更为一般的方程

\begin{equation*}
\frac {d ^ {2} g _ {\lambda} (\xi)}{d \xi^ {2}} + \lambda v (\xi) g _ {\lambda} (\xi) = 0.\tag{7.5}
\end{equation*}

将方程(7.5)的解表示成 $\lambda$ 的幂级数形式:

\begin{equation*}
g _ {\lambda} (\xi) = \sum_ {n = 0} ^ {\infty} \lambda^ {n} g _ {n} (\xi),\tag{7.6}
\end{equation*}

并假定 $\lambda$ 充分小，使级数收敛。利用这个展开式，可得函数 $g_{n}(\xi)$ 和 $g_{n - 1}(\xi)$ 之间的联系：

\begin{equation*}
\frac {d ^ {2} g _ {n} (\xi)}{d \xi^ {2}} + v (\xi) g _ {n - 1} (\xi) = 0.\tag{7.7}
\end{equation*}

在 $n$ = 0 时，方程序列(7.7)的解为

\begin{equation*}
g _ {0} = a _ {0} + b _ {0} \xi ;
\end{equation*}

当 $n = 1$ 时，

\begin{equation*}
\frac {d ^ {2} g _ {1} (\xi)}{d \xi^ {2}} = - v (\xi) g _ {0},
\end{equation*}

因此
\begin{equation*}
g _ {1} (\xi) = - \int_ {0} ^ {\xi} d \xi_ {2} \int_ {0} ^ {\xi_ {2}} d \xi_ {1} v (\xi_ {1}) g _ {0} (\xi_ {1}),
\end{equation*}

当 $n = 2$ 时，

\begin{equation*}
g _ {2} (\xi) = (- 1) ^ {2} \int_ {0} ^ {\xi} d \xi_ {4} \int_ {0} ^ {\xi_ {4}} d \xi_ {3} \int_ {0} ^ {\xi_ {3}} d \xi_ {2} \int_ {0} ^ {\xi_ {2}} d \xi_ {1} v (\xi_ {1}) v (\xi_ {3}) g _ {0} (\xi_ {1});
\end{equation*}

因此，对任意 $n$ 有

\begin{equation*}
\begin{array}{l} g _ {n} (\xi) = (- 1) ^ {n} \int_ {0} ^ {\xi} d \xi_ {2 n} \int_ {0} ^ {\xi_ {2 n}} d \xi_ {2 n - 1} \dots \\ \dots \int_ {0} ^ {\xi_ {2}} d \xi_ {1} [ v (\xi_ {2 n - 1}) v (\xi_ {2 n - 3}) \dots v (\xi_ {1}) ] g _ {0} (\xi_ {1}). \end{array}
\end{equation*}

因为函数 $g_{n}(\xi)$ 在所考虑的区域内解析，所以积分值与积分路径无关。因此，为方便计，所有积分都可沿着矢径 $\xi_{i}=r_{i}e^{i\theta}$ 求积，这里 $\theta$ 值固定。于是

\begin{equation*}
\int_ {0} ^ {\xi_ {i + 1}} \cdot d \xi_ {i} = \int_ {0} ^ {r _ {i + 1}} \cdot e ^ {i \theta} d r _ {i}.
\end{equation*}

为了确定所得的表达式是否是原方程的解，就必须研究级数(7.6)的收敛性，换句话说，需估计值 $g_{n}(\xi)$ 。

由于函数 $v(\xi)$ 和 $g_{0}(\xi)$ 解析，所以成立估计式

\begin{equation*}
\left| v (\xi) \right| <   M, \quad \left| g _ {0} (\xi) \right| <   A,
\end{equation*}

这里 $M$ 和 $A$ 均为常数。因此

\begin{equation*}
\left| g _ {n} (\xi) \right| <   M ^ {n} A \Omega_ {n},
\end{equation*}

这里
\begin{equation*}
\Omega_ {n} = \int_ {0} ^ {r} d r _ {2 n} \dots \int_ {0} ^ {r _ {2}} d r _ {1}.
\end{equation*}

通过直接求积分，不难得

\begin{equation*}
\Omega_ {n} = \int_ {0} ^ {r} d r _ {2 n} \dots \int_ {0} ^ {r _ {2}} d r _ {1} = \frac {r ^ {2 n}}{(2 n) !}.\tag{7.8}
\end{equation*}

现在注意到(7.8)就可写出一连串估值不等式

\begin{equation*}
\left| g _ {n} (\xi) \right| <   \frac {\left| \xi \right| ^ {2 n} M ^ {n} A}{(2 n) !} <   \frac {(a ^ {2} M) ^ {n} A}{(2 n) !},
\end{equation*}

因此对于任意的 $\xi, \lambda$ ，当 $n \to \infty$ 时

\begin{equation*}
\left| \lambda^ {n} g _ {n} (\xi) \right| <   \frac {(a ^ {2} | \lambda | M) ^ {n} A}{(2 n) !} \rightarrow 0.
\end{equation*}

此式本身即已证明了级数(7.6)的收敛性，因为 $a$ 可以取得任意小。级数是一致收敛的，由此推出了函数 $g(\xi)$ 在圆内及边界上，亦即在 $p(\xi)$ 和 $q(\xi)$ 解析的区域内的解析性。再加上当 $n > 0$ 时 $g_{n}(0) = 0$ ，而 $g_0(0) = a_0$ ，因此

\begin{equation*}
g _ {\lambda} (0) = \sum_ {n = 0} ^ {\infty} g _ {n} (0) \lambda^ {n} = a _ {0}
\end{equation*}

和
\begin{equation*}
\frac {d g _ {n}}{d \xi} \Big | _ {\xi = 0} = \left\{ \begin{array}{l l} 0, & \text {当} n > 0 \text {时,} \\ b _ {0}, & \text {当} n = 0 \text {时\text{。}} \end{array} \right.
\end{equation*}

所以
\begin{equation*}
\left. \frac {d g}{d \xi} \right| _ {\xi = 0} = b _ {0}.
\end{equation*}

这样，我们求得的方程(7.4)的解在点 $\xi = 0$ 的邻域内解析，且依赖于 $g(\xi)$ 和 $dg / d\xi$ 在点 $\xi = 0$ 处的值。利用解析开拓的方法可得在除去奇点外整个复平面上的解。

\textbf{定理 7.1}\quad 假定方程 (7.1) 的解 $f_{1}(z)$ 和 $f_{2}(z)$ 是两个在区域 $R$ 内解析的函数，且在某点 $z = z_{0}$ 处

\begin{equation*}
f _ {1} (z _ {0}) = f _ {2} (z _ {0}) \text {和} f _ {1} ^ {\prime} (z _ {0}) = f _ {2} ^ {\prime} (z _ {0}).
\end{equation*}

则对位于 $R$ 内的所有 $z, f_{1}(z) = f_{2}(z)$ .

\textbf{证明}\quad 记

\begin{equation*}
F (z) \equiv f _ {1} (z) - f _ {2} (z).
\end{equation*}

显然
\begin{equation*}
F (z _ {0}) = F ^ {\prime} (z _ {0}) = 0.
\end{equation*}

而由所考虑的微分方程推出 $F''(z_0)=0$ 。将方程微分任意次（由于解析性，这是可能的），容易证明对所有的 $n$ ， $F^{(n)}(z_0) = 0$ 。此意味着函数 $F(z)$ 在 $p(z)$ 和 $q(z)$ 解析的区域内等于零。所以 $f_{1}(z) = f_{2}(z)$ 。定理证毕。

\textbf{例 1}\quad 埃尔米特方程。如果在方程 (7.1) 中置 $p(z) = -2z$ 和 $q(z) = 2\mu$ ，则得方程

\begin{equation*}
f ^ {\prime \prime} (z) - 2 z f ^ {\prime} (z) + 2 \mu f (z) = 0,\tag{7.9}
\end{equation*}

它称为埃尔米特方程。这里 $p(z)$ 和 $q(z)$ 是除点 $z = \infty$ 外（对于 $p(z)$ ）到处解析的函数。如前面所作的那样，可直接求出方程(7.9)形如

\begin{equation*}
f (z) = \sum_ {n = 0} ^ {\infty} a _ {n} z ^ {n}\tag{7.10}
\end{equation*}

的解，这里级数(7.10)处处收敛。将级数（7.10）代入方程(7.9)，就得级数系数之间的关系式：

\begin{equation*}
\sum_ {n} [ (n + 2) (n + 1) a _ {n + 2} - 2 n a _ {n} + 2 \mu a _ {n} ] z ^ {n} = 0,
\end{equation*}

它对所有 $n$ 都成立，因此

\begin{equation*}
(n + 1) (n + 2) a _ {n + 2} + 2 a _ {n} (\mu - n) = 0,
\end{equation*}

或写成
\begin{equation*}
a _ {n + 2} = - \frac {2 (\mu - n)}{(n + 1) (n + 2)} a _ {n}.
\end{equation*}

(7.10)式的表示法根据系数的奇偶性而不同，即

\begin{equation*}
f _ {\text {偶}} = \sum_ {m = 0} ^ {\infty} b _ {m} z ^ {2 m}, \quad b _ {m + 1} = \frac {2 m - \mu}{(m + 1) (2 m + 1)} b _ {m};\tag{7.11a}
\end{equation*}

\begin{equation*}
f _ {\text {奇}} = \sum_ {n = 0} ^ {\infty} c _ {n} z ^ {2 n + 1}, c _ {n + 1} = \frac {2 n + 1 - \mu}{(n + 1) (2 n + 3)} c _ {n}.\tag{7.11b}
\end{equation*}

假定 $\mu$ 是正整数。如果它是偶数，则级数(7.11$a$)将由有限项组成，因当 $m = \frac{\mu}{2}$ 时，系数 $b_{m + 1}$ 等于零，而以后所有项也等于零。类似地，如果 $\mu$ 为奇数，则级数(7.11$b$)为有限项。换句话说，在 $\mu$ 为正整数时，其中有一个解是有限次多项式，这种多项式称为埃尔米特多项式。它有一个奇点，即在点 $z = \infty$ 处有极性。注意，另外一个用无穷级数表示的解在无穷远处有本性奇点。

写出前几个埃尔米特多项式：

\begin{equation*}
\begin{array}{l l} H _ {0} (z) = \text { const }, & H _ {1} (z) = \text { const } z, \\ H _ {2} (z) = \text { const } (1 - z ^ {2}). \end{array}
\end{equation*}

也容易看出，当 $z$ 很大时，级数(7.11)的系数可导致下述形式：

\begin{equation*}
\lim _ {m \to \infty} \frac {b _ {m + 1}}{b _ {m}} = \frac {1}{m} ^ {(1)}, \quad \lim _ {n \to \infty} \frac {c _ {n + 1}}{c _ {n}} = \frac {1}{n}.
\end{equation*}

函数 $e^{z^2}$ 也具有这些性质：

\begin{equation*}
e ^ {z ^ {2}} = \sum_ {p = 0} ^ {\infty} \alpha_ {p} z ^ {2 p} \quad \text {和} \quad \frac {\alpha_ {p + 1}}{\alpha_ {p}} = \frac {1}{p + 1},
\end{equation*}

所以当 $z \to \infty$ 时，函数 $f_{\text{偶}}$ 和 $f_{\text{奇}}$ 的渐近性是一样的：

\begin{equation*}
f _ {\text { 偶 }} \sim e ^ {z ^ {2}} \quad \text { 和 } \quad f _ {\text { 奇 }} \sim e ^ {z ^ {2}},
\end{equation*}

这里记号～表示

\begin{equation*}
\lim _ {z \rightarrow \infty} \left[ \frac {\ln f}{\ln e ^ {z ^ {2}}} \right] = 1.
\end{equation*}

我们求得了函数 $f_{\text{偶}}(z)$ 和 $f_{\text{奇}}(z)$ 的渐近行为。

将变换(7.2)应用于方程(7.9)。在这种情形，

\begin{equation*}
f (z) = e ^ {\frac {1}{2} z ^ {2}} g (z),
\end{equation*}

于是就得方程

\begin{equation*}
\frac {d ^ {2} g}{d z ^ {2}} + \left[ (2 \mu + 1) - z ^ {2} \right] g (z) = 0.
\end{equation*}

如我们已证明了的，它等价于原来的方程。今后我们利用的就是这种形式的埃尔米特方程。

\textbf{例2}\quad 谐振子。一维谐振子的薛定谔方程为

\begin{equation*}
- \frac {h ^ {2}}{2 m} \frac {\partial^ {2}}{\partial x ^ {2}} \psi (x) + \frac {k x ^ {2}}{2} \psi (x) = E \psi (x).
\end{equation*}

这里
\begin{equation*}
\int_ {- \infty} ^ {\infty} | \psi | ^ {2} d x = 1.
\end{equation*}

转到变量 $z = \frac{x}{a}$ ，并采用记号

\begin{equation*}
E = \frac {\hbar}{2} \sqrt {\frac {k}{m}} (2 \mu + 1), \alpha = \frac {\hbar}{\sqrt {k m}}
\end{equation*}

( $\mu$ 是某个整数)。此时方程就变成

\begin{equation*}
\frac {d ^ {2} \psi (z)}{d z ^ {2}} + [ (2 \mu + 1) - z ^ {2} ] \psi (z) = 0,
\end{equation*}

显然，它是埃尔米特方程。所以我们置

\begin{equation*}
\psi = e ^ {- \frac {1}{2} z ^ {2}} f (z),
\end{equation*}

这里 $f(z)$ 满足方程

\begin{equation*}
\frac {d ^ {2} f}{d z ^ {2}} - 2 z \frac {d f}{d z} + 2 \mu f = 0.
\end{equation*}

函数 $f(z)$ 可以是多项式或无穷级数(具有偶次幂或奇次幂)。标准化条件

\begin{equation*}
\int_ {- \infty} ^ {\infty} | \psi (z) | ^ {2} d z = 1
\end{equation*}

给出
\begin{equation*}
\int_ {- \infty} ^ {\infty} e ^ {- z ^ {2}} f ^ {2} (z) d z <   \infty ,\tag{7.12}
\end{equation*}

但无穷级数(7.11)渐近趋向于 $e^{z^{2}}$ 。因为 $\int e^{z^{2}}dz$ 发散，由此推知积分(7.12)无界。所以 $f(z) = H_{\mu}(z)$ （这里 $\mu$ 为正整数），而谐振子的特征值和特征函数为

\begin{equation*}
E = \frac {\hbar}{2} \sqrt {\frac {k}{m}} (2 \mu + 1);
\end{equation*}

\begin{equation*}
\psi (x) = e ^ {- \frac {1}{2} z ^ {2}} H _ {\mu} (z).
\end{equation*}

\section{在正则奇点附近方程的解}

到现在为止，我们建立了下面的结果：所有奇点为已知的方程（7.1）有两个独立解。其中任意一个可用两个常数确定。例如

\begin{equation*}
\begin{array}{r l} & {f _ {1} (z _ {0}) = 0, \quad f _ {1} ^ {\prime} (z _ {0}) = 1 \quad \text {和}} \\ & {f _ {2} (z _ {0}) = 1, \quad f _ {2} ^ {\prime} (z _ {0}) = 0.} \end{array}
\end{equation*}

和通常一样，通解将为 $f_{1}(z)$ 和 $f_{2}(z)$ 的线性组合：

\begin{equation*}
f (z) = a f _ {1} (z) + b f _ {2} (z).
\end{equation*}

函数 $f(z)$ 是方程(7.1)在所有正常点上的解。现在来求在确定类奇点邻域内的解。定义两类奇点。

如果 $z_0$ 是奇点，且在点 $z = z_0$ 处函数 $(z - z_0)p(z)$ 和 $(z - z_0)^2 q(z)$ 解析，即如果 $p(z)$ 有一阶极点，而 $q(z)$ 有二阶极点，则 $z_0$ 称为正则奇点；在其余情形， $z_0$ 就称为非正则奇点。今后我们只考虑正则奇点。

设 $z_0$ 为正则奇点， $z - z_0 = \xi$ 。记

\begin{equation*}
P (\xi) \equiv (z - z _ {0}) p (z) | _ {z - z _ {0} = \xi},
\end{equation*}

\begin{equation*}
Q (\xi) \equiv (z - z _ {0}) ^ {2} q (z) | _ {z - z _ {0} = \xi}.
\end{equation*}

函数 $P(\xi)$ 和 $Q(\xi)$ 在点 $\xi = 0$ 处解析，将这些函数在零点邻域内展开成 $\xi$ 的幂级数：

\begin{equation*}
P (\xi) = \sum_ {m = 0} ^ {\infty} P _ {m} \xi^ {m}, \quad Q (\xi) = \sum_ {m = 0} ^ {\infty} Q _ {m} \xi^ {m}.
\end{equation*}

假定这些级数在 $|\xi| < \xi_0$ 时解析。在 $\xi$ 很小时，可分出方程的奇异部分，即当 $\xi \to 0$ 时趋向于无穷的项：

\begin{equation*}
\frac {d ^ {2} f (z)}{d z ^ {2}} + \frac {P _ {0}}{\xi} \frac {d f (z)}{d z} + \frac {Q _ {0}}{\xi^ {2}} f (z) = 0.\tag{7.13}
\end{equation*}

在求出方程(7.13)的解以后，需求方程(7.1)的解，但已无奇异部分。今试求方程(7.13)形如

\begin{equation*}
f (z) = \xi^ {\sigma}\tag{7.13a}
\end{equation*}

的解，这里 $\sigma$ 待定。为了确定 $\sigma$ ，将解(7.13$a$)代入方程(7.13)，得对于 $\sigma$ 的特征方程

\begin{equation*}
F (\sigma) = \sigma (\sigma - 1) + P _ {0} \sigma + Q _ {0} = 0\tag{7.14}
\end{equation*}

或即
\begin{equation*}
F (\sigma) = \left(\sigma - \alpha_ {1}\right) \left(\sigma - \alpha_ {2}\right) = 0,
\end{equation*}

这里 $\alpha_{1}$ 和 $\alpha_{2}$ 为方程(7.14)的根。

其次，象前面那样将方程(7.1)的解表示成级数形式

\begin{equation*}
f _ {i} (\xi) = \xi^ {\alpha_ {i}} \sum_ {n = 0} ^ {\infty} a _ {n} \xi^ {n} (i = 1, 2).\tag{7.15}
\end{equation*}

从前面知道，系数 $a_0$ 是任意的，因而我们设 $a_0 = 1$ 。容易看出，当 $\xi$ 很小时，象这种形式的解就变成刚才所研究过的方程(7.13)的解。

先考虑 $i = 1$ 时的解(7.15)。对 $f(z)$ 关于 $z$ 微分二次：

\begin{equation*}
\frac {d f}{d z} = \sum_ {n = 0} ^ {\infty} (n + \alpha_ {1}) a _ {n} \xi^ {n + \alpha_ {1} - 1},\tag{7.16}
\end{equation*}

\begin{equation*}
\frac {d ^ {2} f}{d z ^ {2}} = \sum_ {n = 0} ^ {\infty} (n + \alpha_ {1}) (n + \alpha_ {1} - 1) a _ {n} \xi^ {n + \alpha_ {1} - 2}.
\end{equation*}

将方程(7.1)本身表示成

\begin{equation*}
\xi^ {2} \frac {d ^ {2} f (\xi)}{d \xi^ {2}} + \xi P (\xi) \frac {d f (\xi)}{d \xi} + Q (\xi) f (\xi) = 0\tag{7.17}
\end{equation*}

是方便的。将(7.16)式和 $Q(\xi)$ 与 $P(\xi)$ 的展开式代入方程(7.17)，得

\begin{equation*}
\begin{array}{r l} \sum_ {n = 0} ^ {\infty} \xi^ {n + \alpha_ {1}} \left\{a _ {n} [ (n + \alpha_ {1} - 1) (n + \alpha_ {1}) + P _ {0} (n + \alpha_ {1}) + Q _ {0} ] + \right. \\ & \left. + a _ {n - 1} [ (n - 1 + \alpha_ {1}) P _ {1} + Q _ {1} ] + \right. \\ & \left. + a _ {n - 2} [ (n - 2 + \alpha_ {1}) P _ {2} + Q _ {2} ] + \dots \right\} = 0. \end{array}
\end{equation*}

我们依次研究所得级数的每一项。当 $n = 0$ 时

\begin{equation*}
a _ {0} \left[ \alpha_ {1} \left(\alpha_ {1} - 1\right) + P _ {0} \alpha_ {1} + Q _ {0} \right] = a _ {0} F \left(\alpha_ {1}\right) = 0.
\end{equation*}

但 $F(\alpha_{1}) = 0$ ，因为 $\alpha_{1}$ 为特征方程(7.14)的根。因此常数 $a_0$ 可任意选取，如我们所作的那样，假定 $a_0 = 1$。当 $n = 1$ 时，

\begin{equation*}
a _ {1} F \left(\alpha_ {1} + 1\right) + a _ {0} \left[ \alpha_ {1} P _ {1} + Q _ {1} \right] = 0,
\end{equation*}

由此可确定系数 $a_1$ [在 $F(\alpha_1 + 1) \neq 0$ 的条件下]。

一般地，对任意 $n$

\begin{equation*}
a _ {n} F (n + \alpha_ {1}) + \sum_ {m = 1} ^ {n} a _ {n - m} [ (n - m + \alpha_ {1}) P _ {m} + Q _ {m} ] = 0.
\end{equation*}

注意，给定了 $a_0$ ，则在条件 $F(n + \alpha_{1})\neq 0$ 下，由最后一个等式可求出任意一个系数 $a_{n}$ 。如果条件 $F(n + \alpha_{1}) \neq 0$ 对所有 $n$ 满足，则可取（收敛或发散的）无穷级数暂时作为方程(7.17)的形式解。关于收敛性问题，我们将对下面所要考虑的几个特殊情形，作最简单的研究。如果假定对某个 $n$ ， $F(n + \alpha_{1}) = 0$ ，那么这种研究的必要性已失去意义。此时 $n + \alpha_{1}$ 是特征方程的根，换言之， $\alpha_{2} = n + \alpha_{1}$ ，因而在这种条件下差 $\alpha_{2} - \alpha_{1}$ 是整数，级数(7.15)成为多项式。在其余情形，就必须研究收敛性。

情形1. 设 $\alpha_{1} - \alpha_{2}$ 不是整数。因 $P_{m}$ 和 $Q_{m}$ 是解析函数展开成级数时的系数，故

\begin{equation*}
P _ {m} = \frac {1}{2 \pi i} \oint_ {| \xi | = r} \frac {P (\xi)}{\xi^ {m + 1}} d \xi ,
\end{equation*}

\begin{equation*}
Q _ {m} = \frac {1}{2 \pi i} \oint_ {| \xi | = r} \frac {Q (\xi)}{\xi^ {m + 1}} d \xi .
\end{equation*}

如果注意到
\begin{equation*}
\mid P (\xi) \mid \leqslant M, \quad \mid Q (\xi) \mid \leqslant M ^ {\prime},
\end{equation*}

则系数 $P_{m}$ 和 $Q_{m}$ 是容易估计的。 $P_{m}$ 的估计是这样的：

\begin{equation*}
\mid P _ {m} \mid \leqslant \frac {c}{r ^ {m}},
\end{equation*}

这里 $c$ 为常数。类似地，对 $Q_{m}$ 有

\begin{equation*}
\mid Q _ {m} \mid \leqslant \frac {c ^ {\prime}}{r ^ {m}},
\end{equation*}

这里 $c'$ 为常数。显然存在常数 $A$ ，使得

\begin{equation*}
\mid P _ {m} \mid \leqslant \frac {A}{r ^ {m}}, \quad \mid \alpha_ {1} P _ {m} + Q _ {m} \mid \leqslant \frac {A}{r ^ {m}}.\tag{7.18}
\end{equation*}

为了证明级数(7.15)的收敛性[而此即等价于方程(7.1)解存在性的证明]，只需证明系数 $a_{m}$ 满足不等式

\begin{equation*}
\left| a_ {m} \right| \leqslant \left| \frac {k}{r} \right| ^ {m} (k \geqslant 1).\tag{7.19}
\end{equation*}

此时级数(7.15)所表示的函数 $f(\xi)$ (解)解析，而级数本身将有收敛半径$r$/$k$。

现证明不等式(7.19)成立。系数 $a_{0}$ 前面已确定；对于它

\begin{equation*}
\left| a _ {0} \right| \leqslant (k / r) ^ {0}.
\end{equation*}

利用数学归纳法。即假定对所有 $m$ < $n$，有：

\begin{equation*}
\left| a _ {m} \right| \leqslant (k / r) ^ {m},
\end{equation*}

进而再证明此不等式对 $m = n$ 也成立。利用估计式(7.18)来估计和

\begin{equation*}
S = \sum_ {m = 1} ^ {n} a _ {n - m} [ (n - m) P _ {m} + \alpha_ {1} P _ {m} + Q _ {m} ].
\end{equation*}

我们有
\begin{equation*}
S \leqslant \frac {A}{r ^ {n}} \sum_ {m = 1} ^ {n} [ k ^ {n - m} (n - m) + k ^ {n - m} ];
\end{equation*}

又因关系式
\begin{equation*}
\sum_ {m = 1} ^ {n} k ^ {n - m} = 1 + k + k ^ {2} + \dots + k ^ {n - 1} \leqslant n k ^ {n - 1}
\end{equation*}

和
\begin{equation*}
\sum_ {m = 1} ^ {n} \left(k ^ {n - m}\right) (n - m) \leqslant \frac {1}{2} (n - 1) n k ^ {n - 1}
\end{equation*}

成立，故最后估计式可表示如下：

\begin{equation*}
\left| \sum_ {m = 1} ^ {n} a _ {n - m} [ (n - m) P _ {m} + \alpha_ {1} P _ {m} + Q _ {m} ] \right| \leqslant
\end{equation*}

\begin{equation*}
\leqslant \frac {A}{r ^ {n}} k ^ {n - 1} \cdot \frac {(n + 1) n}{2};
\end{equation*}

换句话说，对于系数 $a_{n}$ ，不等式

\begin{equation*}
a _ {n} \leqslant \frac {A k ^ {n - 1} n (n + 1)}{2 F (n + \alpha_ {1}) r ^ {n}} = \frac {A k ^ {n - 1} (n + 1)}{2 (n + \alpha_ {1} - \alpha_ {2})} \frac {1}{r ^ {n}}\tag{7.20}
\end{equation*}

成立。如果选取 $k$ ，使得对所有 $\pmb{n}$

\begin{equation*}
k > \frac {A (n + 1)}{2 \left(n + \alpha_ {1} - \alpha_ {2}\right)},
\end{equation*}

则我们的结论将成立。但表达式

\begin{equation*}
\frac {A (n + 1)}{2 \left(n + \alpha_ {1} - \alpha_ {2}\right)}
\end{equation*}

有不依赖于 $n$ 的极大值，因此存在那样的 $k$，使

\begin{equation*}
\left| a _ {m} \right| \leqslant \left(\frac {k}{r}\right) ^ {m}.\tag{7.21}
\end{equation*}

这就证明了我们的结论，即级数在半径为 $r$/$k$ 的圆内收敛。利用解析延拓的方法可将解推广到此圆之外部。

作出一般结论。

在量 $(\alpha_{1} - \alpha_{2})$ 不等于整数的情形，方程(7.17)有两个用(7.15)式确定的线性无关解，而通解可表示成和

\begin{equation*}
a f _ {1} (z) + b f _ {2} (z)
\end{equation*}

的形式。

对于量 $(\alpha_{1} - \alpha_{2})$ 等于整数时的情形，目前还只知道下面这些：

如果 $(\alpha_{1} - \alpha_{2})$ 等于零，则方程有一个解 $f_{1} = f_{2}$。如果 $(\alpha_{1} - \alpha_{2})$ 是正整数，则 $f_{2}(z)$ 存在，而 $f_{1}$ 未知；如果 $(\alpha_{1} - \alpha_{2})$ 是负整数，则 $f_{1}(z)$ 存在，而 $f_{2}$ 未知。

情形2. 差 $(\alpha_{1} - \alpha_{2})$ 等于正整数。假定还存在一个除 $f_{1}(z)$ 以外的解。它可以从 $f_{2}(z) = f_{1}(\xi)g(\xi)$ 形式中去寻找。于是

\begin{equation*}
\begin{array}{r l} f _ {2} ^ {\prime} (\xi) & = f _ {1} ^ {\prime} (\xi) g (\xi) + f _ {1} (\xi) g ^ {\prime} (\xi), \\ f _ {2} ^ {\prime \prime} (\xi) & = f _ {1} ^ {\prime \prime} (\xi) g (\xi) + 2 f _ {1} ^ {\prime} (\xi) g ^ {\prime} (\xi) + \\ & \quad + f _ {1} (\xi) g ^ {\prime \prime} (\xi). \end{array}\tag{7.22}
\end{equation*}

将(7.22)代入原方程后得

\begin{equation*}
\begin{array}{r l} & \xi^ {2} f _ {1} g ^ {\prime \prime} (\xi) + 2 \xi^ {2} f _ {1} ^ {\prime} g ^ {\prime} (\xi) + \xi P (\xi) f _ {1} g ^ {\prime} (\xi) + \\ & \quad + g [ \xi^ {2} f _ {1} ^ {\prime \prime} + \xi P f _ {1} ^ {\prime} + Q f _ {1} ] = 0, \end{array}
\end{equation*}

因括号内的项等于零，故

\begin{equation*}
\xi^ {2} f _ {1} \frac {d ^ {2} g}{d \xi^ {2}} + 2 \xi^ {2} \frac {d f _ {1}}{d \xi} \frac {d g}{d \xi} + \xi P f _ {1} \frac {d g}{d \xi} = 0.
\end{equation*}

这样，对函数 $g'(\xi)$ 有下面的一阶方程

\begin{equation*}
\left(\frac {d g}{d \xi}\right) ^ {- 1} \cdot \frac {d}{d \xi} \left(\frac {d g}{d \xi}\right) = - \frac {2}{f _ {1}} \frac {d f _ {1}}{d \xi} - \frac {P}{\xi}.\tag{7.23}
\end{equation*}

它的积分由下式给出：

\begin{equation*}
\ln \frac {d g}{d \xi} = - 2 \ln f _ {1} - P _ {0} \ln \xi - \sum_ {m = 1} ^ {\infty} P _ {m} \xi^ {m} \frac {1}{m} + \text { const }.\tag{7.24}
\end{equation*}

这里利用了表达式 $P(\xi) = \sum_{m=0}^{\infty} P_m \xi^m$ 。因此

\begin{equation*}
\frac {d g}{d \xi} = \frac {A \exp \left(- \sum_ {m = 1} ^ {\infty} \frac {P _ {m}}{m} \xi^ {m}\right)}{f _ {1} ^ {2} \xi^ {P _ {0}}} = \frac {A \varphi (\xi)}{\xi^ {2 a _ {1} + P _ {0}}},\tag{7.25}
\end{equation*}

这里
\begin{equation*}
\varphi (\xi) = \frac {\exp \left[ - \sum_ {m = 1} ^ {\infty} \frac {P _ {m}}{m} \xi^ {m} \right]}{\left(\sum_ {m = 0} ^ {\infty} a _ {m} \xi^ {m}\right) ^ {2}},
\end{equation*}

因为 $f_{1} = \xi^{\alpha_{1}}\sum_{m = 0}^{\infty}a_{m}\xi^{m}$ 。显然函数 $\varphi (\xi)$ 在点 $\xi = 0$ 处解析，因此它可表示成级数

\begin{equation*}
\varphi (\xi) = \sum_ {m = 0} ^ {\infty} c _ {m} \xi^ {m},
\end{equation*}

这里 $c_{0} = 1$ ，因为 $a_0 = 1$ 。注意到等式

\begin{equation*}
F (\alpha) = \left(\alpha - \alpha_ {1}\right) \left(\alpha - \alpha_ {2}\right) = \alpha (\alpha - 1) + P _ {0} \alpha + Q _ {0} = 0,
\end{equation*}

得
\begin{equation*}
P _ {0} + 2 \alpha_ {1} = \alpha_ {1} - \alpha_ {2} + 1 = S + 1.
\end{equation*}

这里 $S$ 是正整数。所以(7.25)可改写成

\begin{equation*}
\begin{array}{r l} \frac {d g}{d \xi} & = A \frac {\sum_ {m = 0} ^ {\infty} c _ {m} \xi^ {m}}{\xi^ {S + 1}} = \\ & = A \left[ \frac {c _ {0}}{\xi^ {S + 1}} + \frac {c _ {1}}{\xi^ {S}} + \dots + \frac {c _ {S}}{\xi} + c _ {S + 1} + \xi c _ {S + 2} + \dots \right] \end{array}\tag{7.26}
\end{equation*}

此式积分后就得函数 $g(\xi)$ 的表达式

\begin{equation*}
g (\xi) = A c _ {s} \ln \xi + \xi^ {- s} \sum_ {m = 0} ^ {\infty} d _ {m} \xi^ {m} + \text { const },\tag{7.27}
\end{equation*}

这里函数 $\sum d_{m}\xi^{m}$ 解析，而常数等于零，因为我们选取 $f_{2} = f_{1}g$ ，且 $cf_{1}$ 型的解对我们是没有意义的[因为我们要求的是与 $f_{1}(z)$ 线性无关的解]。

这样
\begin{equation*}
f _ {2} (\xi) = A c _ {s} f _ {1} (\xi) \ln \xi + \xi^ {\alpha_ {2}} \sum_ {m = 0} ^ {\infty} b _ {m} \xi^ {m}\tag{7.28}
\end{equation*}

(级数 $\sum b_{m}z^{m}$ 是解析函数 $f_{1}$ 和 $g$ 互乘的结果)。正如我们看到的，在 $S$=0 时，因为 $c_{0}=1$ ，函数 $f_{2}(\xi)$ 有 $a\ln z$ 型的项；如果 $S\neq0$ ，则 $c_{s}$ 可等于零。大家知道，函数 $\ln z$ 在复平面内有支点 $z$=0（这个问题在第六章中已讨论过）。

\section{贝塞耳方程}

下面形状的二阶微分方程称为贝塞尔方程：

\begin{equation*}
\frac {d ^ {2} f}{d z ^ {2}} + \frac {1}{z} \frac {d f}{d z} + \left(1 - \frac {n ^ {2}}{z ^ {2}}\right) f = 0.\tag{7.29}
\end{equation*}

显然， $z = 0$ 是正则奇点。因为函数 $p = 1 / z$ 和 $q = 1 - (n^2 / z^2)$ 满足前节所讲的条件。根据公式(7.15)，有级数形式的解：

\begin{equation*}
f (z) = z ^ {a} \sum_ {m = 0} ^ {\infty} a _ {m} z ^ {m},\tag{7.29a}
\end{equation*}

这里 $a_0 \neq 0$ 。将 $f(z)$ 代入方程后，就得确定级数系数的关系式

\begin{equation*}
a _ {m} [ (\alpha + m) (\alpha + m - 1) + (\alpha + m) - n ^ {2} ] + a _ {m - 2} = 0.\tag{7.30}
\end{equation*}

当 $m = 0$ 时，(7.30)就转变为特征方程

\begin{equation*}
\alpha (\alpha - 1) + \alpha - n ^ {2} = 0.
\end{equation*}

由此推出 $\alpha = \pm n$ ; 因而可能有两个解。

当 $m = 1$ 时，系数 $a_1 = 0$ ，因此所有奇次幂系数等于零。幂级数(7.29)仅含有 $z$ 的偶次幂项，并且具有下面的形式

\begin{equation*}
f _ {\pm} (z) = z ^ {\pm n} \sum_ {l = 0} ^ {\infty} b _ {l} z ^ {2 l}, \quad b _ {l} = - \frac {b _ {l - 1}}{4 (l + \alpha) l}.\tag{7.31}
\end{equation*}

如果 $n$ 不是整数，则(7.31)中关于 $b_{1}$ 表达式的分母不等于零。因为(7.31)带有正负号的解线性无关，故通解将是解 $f_{1}$ 和 $f_{2}$ 的线性组合。

如果 $n$ 是整数，则函数

\begin{equation*}
f _ {1} (z) = z ^ {| n |} \sum_ {l = 0} ^ {\infty} b _ {l} z ^ {2 l},\tag{7.32}
\end{equation*}

\begin{equation*}
f _ {2} (z) = c _ {S} f _ {1} (z) \ln z + z ^ {- | n |} \sum_ {l = 0} ^ {\infty} c _ {l} z ^ {2 l}\tag{7.33}
\end{equation*}

可作为解。这两个函数称为贝塞耳函数。下面我们将研究贝塞耳函数的一些基本性质。

积分表示法是解微分方程的强有力的方法。它在我们这个特殊情形，给出了用复平面上的积分来定义贝塞耳函数的可能性。定义 $J_{n}(z)$ 为

\begin{equation*}
J _ {n} (z) \equiv \frac {1}{2 \pi i} \left(\frac {z}{2}\right) ^ {n} \int_ {C} \frac {e ^ {(t - \frac {z ^ {2}}{4 t})}}{t ^ {n + 1}} d t,\tag{7.34}
\end{equation*}

这里在复平面内的围道如图40所示。沿着负实半轴作割缝，并将此定义推广到非整数 $n$ 。被积函数除点 $t = 0$ 外到处解析。因而积分可沿着绕过点 $t = 0$ 的任意道路。注意当 $n$ 为整数时，点 $t = 0$ 为极点；而当 $n$ 为分数时，为支点。为了选择单值分支，在割缝上假定

\begin{figure}[htbp]
\centering
\includegraphics[width=0.55\textwidth]{fig40.pdf}
\par\vspace{2pt}
{\small 图 40.}
\end{figure}

\begin{equation*}
t ^ {n} = | t | ^ {n} e ^ {i n \pi}.
\end{equation*}

今证明，这样确定的函数 $J_{n}(z)$ 满足贝塞耳方程。假定

\begin{equation*}
\Phi_ {n} \equiv z ^ {n} e ^ {- z ^ {2} / 4 t},
\end{equation*}

\begin{equation*}
\begin{array}{l} \Phi_ {n} ^ {\prime} = \left[ n z ^ {n - 1} - \frac {z ^ {n + 1}}{2 t} \right] e ^ {- z ^ {2} / 4 t}, \\ \Phi_ {n} ^ {\prime \prime} = \left[ n (n - 1) z ^ {n - 2} - \frac {1}{2 t} (2 n + 1) z ^ {n} + \frac {z ^ {n + 2}}{4 t ^ {2}} \right] e ^ {- z ^ {2} / 4 t}. \end{array} \tag{7.35}
\end{equation*}

将(7.35)代入贝塞耳方程得

\begin{equation*}
\begin{array}{r l} & \Phi_ {n} ^ {\prime \prime} + \frac {1}{z} \Phi_ {n} ^ {\prime} + \left(1 - \frac {n ^ {2}}{z ^ {2}}\right) \Phi_ {n} = \\ & = e ^ {- z ^ {2} / 4 t} \left[ z ^ {n} \left(1 - \frac {n + 1}{t}\right) + \frac {z ^ {n + 2}}{4 t ^ {2}} \right], \end{array}
\end{equation*}

即
\begin{equation*}
\begin{array}{l} \frac {e ^ {t}}{t ^ {n + 1}} \left[ \Phi_ {n} ^ {\prime \prime} + \frac {1}{z} \Phi_ {n} ^ {\prime} + \left(1 - \frac {n ^ {2}}{z ^ {2}}\right) \Phi_ {n} \right] = \\ = \frac {d}{d t} \left[ \frac {e ^ {(t - \frac {z ^ {2}}{4 t})}}{t ^ {n + 1}} \right] z ^ {n}. \end{array}
\end{equation*}

与 $z$ 有关的部分全部包含在函数 $\Phi_{n}$ 之内。由此直接推出等式

\begin{equation*}
\begin{array}{l} J _ {n} ^ {\prime \prime} (z) + \frac {1}{z} J _ {n} ^ {\prime} (z) + \left(1 - \frac {n ^ {2}}{z ^ {2}}\right) J _ {n} (z) = \\ = \frac {1}{2 \pi i} \int_ {C} \left(\frac {z}{2}\right) ^ {n} \frac {d}{d t} \frac {e ^ {(t - \frac {z ^ {2}}{4 t})}}{t ^ {n + 1}} d t. \end{array}
\end{equation*}

但位于右边的积分等于零，因被积函数当 $t \to -\infty - iO_{+}$ 和 $t \to -\infty + iO_{+}$ 时趋向于零。所以 $J_{n}(z)$ 是贝塞耳函数。今变换积分(7.34)。注意到 $t$ 在任何时刻都不为零，将被积函数展开成 $t$ 的幂级数，则 $J_{n}(z)$ 变成

\begin{equation*}
J _ {n} (z) = \frac {1}{2 \pi i} \left(\frac {z}{2}\right) ^ {n} \int_ {C} \frac {e ^ {t}}{t ^ {n + 1}} \left[ 1 - \frac {z ^ {2}}{4 t} + \frac {z ^ {4}}{2 ! 1 6 t ^ {2}} - \dots \right] d t.\tag{7.36}
\end{equation*}

定义函数 $\Gamma(l)$ 为

\begin{equation*}
\frac {1}{\Gamma (l)} \equiv \frac {1}{2 \pi i} \int_ {c} \frac {e ^ {t}}{t ^ {l}} d t.\tag{7.37}
\end{equation*}

函数 $\Gamma(l)$ 在数学中经常碰到，并称之为 $\Gamma$ 函数(伽玛函数)。利用 $\Gamma(l)$ ，展开式(7.36)可用下面的方式表示：

\begin{equation*}
J _ {n} (z) = z ^ {n} \sum_ {m = 0} ^ {\infty} \frac {(- z ^ {2}) ^ {m}}{\Gamma (m + n + 1) 2 ^ {n + 2 m} m !},
\end{equation*}

它单值地确定 $J_{n}(z)$ 。

再建立 $\Gamma$ 函数的几个性质。

\begin{equation*}
\frac {1}{\Gamma (l)} = \frac {1}{2 \pi i} \int_ {c} \frac {e ^ {t}}{t ^ {l}} d t =
\end{equation*}

\begin{equation*}
\begin{array}{l} = \frac {1}{2 \pi i} \left[ - \frac {e ^ {t}}{(l - 1) t ^ {l - 1}} \Bigg | _ {- \infty - \mathrm{i} \circ_ {+}} ^ {- \infty + \mathrm{i} \circ_ {+}} + \frac {1}{l - 1} \int_ {C} \frac {e ^ {t}}{t ^ {l - 1}} d t \right] = \\ = \frac {1}{l - 1} \frac {1}{\Gamma (l - 1)}. \end{array}
\end{equation*}

当 $l$ 为整数时，积分可一直继续下去而得

\begin{equation*}
\Gamma (l) = (l - 1) (l - 2) \dots 2 \Gamma (1),
\end{equation*}

这里显然有
\begin{equation*}
[ \Gamma (1) ] ^ {- 1} = \frac {1}{2 \pi i} \int_ {C} \frac {e ^ {t}}{t} d t,
\end{equation*}

而函数 $e^t / t$ 除 $t = 0$ 外处处解析。显然不必作割缝，我们将沿着绕原点的圆周积分。初等积分法给出

\begin{equation*}
\Gamma (1) = 1, \text {即} \quad \Gamma (l) = (l - 1)!.
\end{equation*}

对于非整数值 $l$ 的 $\Gamma$ 函数，可看作是阶乘的推广。如果 $l$ 为零或负数，则函数 $[\Gamma(l)]^{-1}$ 是解析的。所以积分(7.37)等于零。换言之

\begin{equation*}
\Gamma (0) = \Gamma (- 1) = \Gamma (- 2) = \dots = \infty .
\end{equation*}

现在回到贝塞耳函数的研究。当 $n$ 为非整数时，得线性无关的两个解 $J_{n}$ 和 $J_{-n}$ 。如果 $n$ 是正整数，则 $\Gamma$ 函数仅当 $m - n + 1 \geqslant 1$ （当 $m \geqslant n$ ）时才不为无穷。因而在 $J_{-n}$ 的展开式中对于 $m < n$ 的项应取零值。最后，在 $n$ 为整数时，贝塞耳方程的通解具有下面的形式：

\begin{equation*}
\begin{array}{l} f (z) = A J _ {| n |} + B Y _ {| n |}, \\ Y _ {| n |} (z) = J _ {| n |} (z) \ln z + z ^ {- | n |} \sum_ {m = 0} ^ {\infty} c _ {m} z ^ {2 m}; \end{array}\tag{7.38}
\end{equation*}

$J_{|n|}$ ——第一类贝塞耳函数，

$Y_{|n|}(z)$ ——第二类贝塞耳函数。

利用贝塞耳函数的积分表达式，作出几个递推公式。根据积分表达式，

\begin{equation*}
- \varphi_ {n} ^ {\prime} (z) + \frac {n}{z} \varphi_ {n} (z) = \frac {z ^ {n + 1}}{2 t} e ^ {- z ^ {2} / 4 t},
\end{equation*}

或
\begin{equation*}
\frac {1}{2 ^ {n}} \frac {e ^ {t}}{t ^ {n + 1}} \left[ - \varphi_ {n} ^ {\prime} (z) + \frac {n}{z} \varphi_ {n} (z) \right] = \frac {z ^ {n + 1}}{2 ^ {n + 1} t ^ {n + 2}} e ^ {(t - \frac {z ^ {2}}{4 t})}.
\end{equation*}

将最后的等式沿着图40所示的围道积分，得下面的关系式

\begin{equation*}
- \frac {d J _ {n}}{d z} + \frac {n}{z} J _ {n} (z) = J _ {n + 1} (z).\tag{7.39}
\end{equation*}

类似地可证明

\begin{equation*}
J _ {n + 1} (z) + J _ {n - 1} (z) = \frac {2 n}{z} J _ {n} (z) \text{。}\tag{7.40}
\end{equation*}

特别有趣的是半奇阶的贝塞耳函数。先讲函数 $J_{1/2}(z)$ 。为了计算这个函数展开式的系数，必须知道函数 $\Gamma\left(m + \frac{3}{2}\right)$ 。它等于

\begin{equation*}
\begin{array}{r l} \Gamma \left(m + \frac {3}{2}\right) & = \left(m + \frac {1}{2}\right) \left(m - \frac {1}{2}\right) \dots \Gamma \left(\frac {1}{2}\right) = \\ & = \frac {(2 m + 1) !}{m ! 2 ^ {2 m + 1}} \Gamma \left(\frac {1}{2}\right). \end{array} \tag{7.41}
\end{equation*}

现在可得关于 $J_{1/2}(z)$ 的十分简单的公式：

\begin{equation*}
J _ {1 / 2} (z) = z ^ {1 / 2} \sum_ {m = 0} ^ {\infty} \frac {(- z ^ {2}) ^ {m} m ! 2 ^ {2 m + 1}}{\Gamma \left(\frac {1}{2}\right) (2 m + 1) ! 2 ^ {2 m + 1 / 2} m !} =
\end{equation*}

\begin{equation*}
\begin{array}{l} = \left(\frac {2}{z}\right) ^ {1 / 2} \sum_ {m = 0} ^ {\infty} \frac {(- 1) ^ {m} z ^ {2 m + 1}}{\Gamma \left(\frac {1}{2}\right) (2 m + 1) !} = \\ = \left(\frac {2}{z}\right) ^ {1 / 2} \frac {\sin z}{\Gamma \left(\frac {1}{2}\right)}. \end{array}
\end{equation*}

由定义
\begin{equation*}
\frac {1}{\Gamma \left(\frac {1}{2}\right)} \equiv \frac {1}{2 \pi i} \int_ {c} \frac {e ^ {t}}{t ^ {1 / 2}} d t.
\end{equation*}

\begin{figure}[htbp]
\centering
\includegraphics[width=0.42\textwidth]{fig41.pdf}
\par\vspace{2pt}
{\small 图 41.}
\end{figure}

将积分围道分成三部分(图41)，在围道 $C_{1}$ 上

\begin{equation*}
t = \xi e ^ {i \theta} \quad \text {和} \quad d t = i \xi e ^ {i \theta} d \theta ,
\end{equation*}

半径 $\xi$ 可趋向于零，围道 $C_{1}$ 缩成一点；因此

\begin{equation*}
\begin{array}{l} \left| t \right| = \xi ; \\ \frac {1}{\left| t ^ {1 / 2} \right|} = \frac {1}{\sqrt {\xi}}, \end{array}
\end{equation*}

从而当 $\xi \to 0$ 时

\begin{equation*}
\left| \int_ {C _ {1}} \frac {e ^ {t}}{t ^ {1 / 2}} d t \right| <   \frac {2 \pi \cdot \xi \cdot e ^ {\xi}}{\sqrt {\xi}} \textcircled {2} \rightarrow 0.
\end{equation*}

在围道 $C_2$ 上， $t = x + iO_{+}$ ，且由定义（我们已选取了相应的

\textcircled{1} 此处俄译本误为： $|t|=\xi\frac{1}{|t^{1/2}|}\approx\frac{1}{\sqrt{\xi}}$ 。——校注

\textcircled{2} 此处俄译本为： $\frac{2\pi\xi}{\sqrt{\xi}}\rightarrow 0.$分支)

\begin{equation*}
t ^ {1 / 2} = | x | ^ {1 / 2} e ^ {i \frac {\pi}{2}} = i | x | ^ {1 / 2}.
\end{equation*}

所以
\begin{equation*}
\frac {1}{\Gamma \left(\frac {1}{2}\right)} = - \frac {1}{2 \pi i} \int_ {0} ^ {\infty} \frac {e ^ {- | x |}}{i | x | ^ {1 / 2}} d x + \frac {1}{2 \pi i} \int_ {C _ {3}} \frac {e ^ {t}}{t ^ {1 / 2}} d t.
\end{equation*}

但沿 $C_3$ 的积分和沿 $C_2$ 的积分是一样的。因而

\begin{equation*}
\frac {1}{\Gamma \left(\frac {1}{2}\right)} = - \frac {2}{2 \pi i} \int_ {0} ^ {\infty} \frac {e ^ {- | x |}}{i | x | ^ {1 / 2}} d | x |.
\end{equation*}

借助于变换 $y \to (x)^{1/2}$ ，这个积分可直接计算：

\begin{equation*}
\frac {1}{\Gamma \left(\frac {1}{2}\right)} = \frac {2}{\pi} \int_ {0} ^ {\infty} e ^ {- y ^ {2}} d y = \frac {1}{\sqrt {\pi}}.
\end{equation*}

换言之，函数 $J_{1 / 2}(z)$ 可用初等函数来表示：

\begin{equation*}
J _ {1 / 2} (z) = \left(\frac {2}{\pi z}\right) ^ {1 / 2} \sin z.\tag{7.42}
\end{equation*}

现在回到贝塞耳微分方程(7.29)，并作变换

\begin{equation*}
f = \frac {\varphi (z)}{\sqrt {z}};
\end{equation*}

则对函数 $\varphi$ 有

\begin{equation*}
\frac {d ^ {2} \varphi (z)}{d z ^ {2}} + \left(1 - \frac {n ^ {2} - \frac {1}{4}}{z ^ {2}}\right) \varphi (z) = 0.\tag{7.43}
\end{equation*}

注意，当 $n = \frac{1}{2}$ 时，它的解之一为

\begin{equation*}
\varphi_ {1 / 2} (z) = \sin z = \sqrt {\frac {\pi z}{2}} J _ {1 / 2} (z),\tag{7.44}
\end{equation*}

而对于下面的半奇阶的函数

\begin{equation*}
\varphi_ {n} (z) = \sqrt {\frac {\pi z}{2}} J _ {n} (z)
\end{equation*}

也可计算。例如，根据递推公式

\begin{equation*}
\varphi_ {1 / 2} = \sin z, \quad \varphi_ {3 / 2} = \frac {\sin z}{z} - \cos z,
\end{equation*}

\begin{equation*}
\varphi_ {- 1 / 2} = \cos z, \quad \varphi_ {- 3 / 2} = - \left(\frac {\cos z}{z} + \sin z\right).
\end{equation*}

附注：半奇阶贝塞耳函数可用三角函数和 $z$ 的不同幂次来表示。函数 $\varphi(z)$ 常称为球贝塞耳函数。

\textbf{例 3}\quad 球的振动。大家知道，描写半径为 $a$ 的球面振动微分方程为

\begin{equation*}
- \Delta \psi = \lambda \psi .\tag{7.45}
\end{equation*}

附注：方程(7.45)是在球面坐标系中拉普拉斯算子特征函数的方程，且有边界条件 $\psi = 0$ （在球面上）。

由于球的对称性，可寻求形如

\begin{equation*}
\psi (r, \theta , \varphi) = \sum_ {l, m} R _ {l} (r) Y _ {l m} (\theta , \varphi)\tag{7.46}
\end{equation*}

的解。这里 $Y_{l,m}(\theta, \varphi)$ 和 $R_l(r)$ 满足相应的方程

\begin{equation*}
L ^ {2} Y _ {l m} (\theta , \varphi) = l (l + 1) Y _ {l m} (\theta , \varphi),\tag{7.47}
\end{equation*}

和
\begin{equation*}
- \frac {1}{r ^ {2}} \frac {\partial}{\partial r} \left(r ^ {2} \frac {d R _ {l}}{d r}\right) + \frac {l (l + 1)}{r ^ {2}} R _ {l} = \lambda R _ {l}.\tag{7.48}
\end{equation*}

在方程(7.48)中作变量代换 $R_{i}\rightarrow M_{i} / r$ ，则

\begin{equation*}
- \frac {\partial^ {2} M _ {l}}{\partial r ^ {2}} + \left[ \frac {l (l + 1)}{r ^ {2}} - \lambda \right] M _ {i} = 0,
\end{equation*}

而如果 $z = r \sqrt{\lambda}$ 和 $l + \frac{1}{2} = n$ ，则我们得

\begin{equation*}
\frac {d ^ {2} M _ {l}}{d z ^ {2}} + \left(1 - \frac {n ^ {2} - \frac {1}{4}}{z ^ {2}}\right) M _ {l} = 0,\tag{7.49}
\end{equation*}

因为
\begin{equation*}
l (l + 1) = n ^ {2} - \frac {1}{4}. ^ {(1)}
\end{equation*}

如果 $l$ 是整数，则 $n$ 是奇数之半。这样一来，方程(7.49)的通解是

\begin{equation*}
M _ {l} = A _ {l} \varphi_ {l + 1 / 2} (z) + B _ {l} \varphi_ {- (l + 1 / 2)} (z),
\end{equation*}

这里 $\varphi_{n}(z)$ 是球面贝塞耳函数。

现在阐明加在解上的某种限制的物理要求是怎样的。要求当 $r \to 0$ 时使 $\psi(r, \theta, \varphi)$ 保持有限。因 $R_{l}(r) = M_{l} / r$ ，所以当 $r \to 0$ 时， $M_{l}$ 必须比 $r$ 更快地趋向于零。也注意对 $l = 0$

\begin{equation*}
\lim _ {r \rightarrow 0} \varphi_ {1 / 2} (r) = 0, \quad \lim _ {r \rightarrow 0} \varphi_ {- 1 / 2} (r) = 1,
\end{equation*}

而当 $l > 0$ 时

\begin{equation*}
\lim _ {r \rightarrow 0} \varphi_ {l + 1 / 2} (r) = 0, \quad \lim _ {r \rightarrow 0} \varphi_ {- (l + 1 / 2)} (r) = \infty .
\end{equation*}

因此，仅是下标为正值的解才有物理意义：

\begin{equation*}
\psi (r, \theta , \varphi) = \sum_ {l, m} \frac {A _ {l} \varphi_ {l + 1 / 2} (r \sqrt {\lambda})}{r} Y _ {l m} (\theta , \varphi).\tag{7.50}
\end{equation*}

此外，还必须满足边界条件 $\psi(a, \theta, \varphi) = 0$ 。因 $Y_{lm}(\theta, \varphi)$ 在球面上不等于零，故要求使 $\varphi_{l+1/2}(a\sqrt{\lambda}) = 0$ 。这个条件是关于 $\lambda$ 的方程。它的根给出了拉普拉斯算子的特征值（即球振动特征频率）。在和式(7.50)中的每个 $l$ 都对应了一个函数 $\varphi$ ，而函数 $\psi$ 仅当除一个以外的所有 $A_l$ 都等于零时，才在 $r = a$ 时恒等于零（对所有 $\theta$ 和 $\varphi$ ）。

特征函数的最终表达式为

\begin{equation*}
\psi_ {l m n} = \frac {\varphi_ {l + 1 / 2} (\sqrt {\lambda_ {n l}} r)}{r} A _ {l n} Y _ {l m} (\theta , \varphi).\tag{7.51}
\end{equation*}

\textbf{例 4}\quad 圆柱体振动。方程和球的一样:

\begin{equation*}
- \Delta \psi = \lambda \psi ,\tag{7.52}
\end{equation*}

但这是在柱面坐标系中, 因为边界条件 $\psi = 0$ 是给在柱面上, 即 $\rho = a$ .

显然，问题的解可表示为

\begin{equation*}
\psi (\rho , \theta , z) = F _ {k m} (\rho) e ^ {i m \varphi} e ^ {i k z}.
\end{equation*}

如果注意到函数 $\exp (im\varphi)$ 是拉普拉斯算子在圆上的特征函数，而 $\exp (ikz)$ 是在 $z$ 轴上的特征函数，这种表示法就显得是合理的。

函数 $F_{km}(\rho)$ 满足方程

\begin{equation*}
\begin{array}{r l} & \frac {d ^ {2} F _ {k m} (\rho)}{d \rho^ {2}} + \frac {1}{\rho} \frac {d F _ {k m} (\rho)}{d \rho} + \\ & \quad + \left[ (\lambda - k ^ {2}) - \frac {m ^ {2}}{\rho^ {2}} \right] F _ {k m} (\rho) = 0. \end{array}\tag{7.53}
\end{equation*}

如果置 $z = \rho \sqrt{\lambda - k^2}$ ，则

\begin{equation*}
\frac {d ^ {2} F}{d z ^ {2}} + \frac {1}{z} \frac {d F}{d z} + \left(1 - \frac {m ^ {2}}{z ^ {2}}\right) F = 0,\tag{7.54}
\end{equation*}

这里 $m$ 是整数。此方程的通解是

\begin{equation*}
F (\rho) = A J _ {| m |} (z) + B Y _ {| m |} (z).\tag{7.55}
\end{equation*}

由当 $\rho\rightarrow0$ 时函数有界性的要求推出 $B$=0，或

\begin{equation*}
F _ {k m} (\rho) = A J _ {| m |} \left(\sqrt {\lambda - k ^ {2}} \rho\right),
\end{equation*}

同时
\begin{equation*}
J _ {| m |} (\sqrt {\lambda - k ^ {2}} a) = 0.
\end{equation*}

(后面的要求由有界性条件推出，且是关于 $\lambda$ 的方程)，它的根是我们问题的特征值，依赖于三个下标。这样

\begin{equation*}
\psi_ {n k m} = J _ {| m |} (\sqrt {\lambda_ {n} - k ^ {2}} \rho) e ^ {i m \varphi} e ^ {i k z}.\tag{7.56}
\end{equation*}
